%========================================
% Plantilla artículo - Humanidades Digitales / NLP
% Compilar con: pdfLaTeX + BibTeX + pdfLaTeX + pdfLaTeX
%========================================
\documentclass[11pt,a4paper]{article}

% --- Codificación y tipografía ---
\usepackage[utf8]{inputenc}   % Si usas LuaLaTeX/XeLaTeX, quita inputenc y usa fontspec
\usepackage[T1]{fontenc}
\usepackage{lmodern}
\usepackage{microtype}

% --- Idioma ---
\usepackage[spanish,es-tabla]{babel}  % "Tabla" en lugar de "Cuadro", etc.
\usepackage{csquotes}

% --- Diseño de página ---
\usepackage[margin=1in]{geometry}
\usepackage{setspace}
%\onehalfspacing  % Activa si quieres interlineado 1.5

% --- Matemáticas y figuras/tablas ---
\usepackage{amsmath,amssymb}
\usepackage{graphicx}
\usepackage{booktabs}
\usepackage{tabularx}
\usepackage{siunitx}

% --- Hipervínculos y referencias cruzadas ---
\usepackage[hidelinks,colorlinks=true,linkcolor=blue,citecolor=blue,urlcolor=blue]{hyperref}
\usepackage[nameinlink,capitalise,noabbrev]{cleveref}

% --- Citas/Bibliografía (estilo social sciences) ---
\usepackage[round,authoryear]{natbib} % \citep, \citet
\bibliographystyle{apalike}           % Elige: apalike / plainnat / chicago, etc.

% --- Comandos útiles ---
\newcommand{\keywords}[1]{\par\noindent\textbf{Palabras clave —} #1}
\newcommand{\JEL}[1]{\par\noindent\textbf{Clasificación —} #1}

% --- Metadatos ---
\title{Ciencia y técnica en la esfera pública de Colombia: \\
	Análisis del discurso en columnas de opinión con metodologías de humanidades digitales}

\author{
	Karen Melissa Gomez Montoya\\
	Biviana Marcela Suárez Sierra\\
	\small Universidad EAFIT, Medellín, Colombia\\
	\small \texttt{tu-email@eafit.edu.co}
}

\date{\today}

%========================================
\begin{document}
	\maketitle
	
	\begin{abstract}
		Es labor de las humanidades y ciencias sociales entender cómo se utiliza el lenguaje y cómo este refleja y genera cambios en la sociedad, en particular, cómo circulan y se transforman los discursos científicos. Este trabajo integra metodologías de ciencias de datos aplicadas a las humanidades digitales para analizar el lugar que ocupan los discursos científicos y técnicos en la opinión pública colombiana. Con un corpus de $\sim$8{,}000 columnas (2018--2020), aplicamos un pipeline de PLN (tokenización, lematización), modelado de tópicos y análisis de coocurrencias, junto con una evaluación cualitativa contextual, para responder: (i) ¿cuáles son los usos del discurso científico en la esfera pública de Colombia?, y (ii) ¿cómo se valora socialmente dicho discurso?
	\end{abstract}
	
	\keywords{Humanidades digitales; Análisis del discurso; Procesamiento del lenguaje natural; Opinión pública; Colombia}
	%\JEL{(si aplica)}
	
	1. Introducción breve del hallazgo
	
	Contextualiza qué parte del objetivo aborda el hallazgo (ej. frecuencia de términos, tópicos emergentes, valoración discursiva).
	
	Usa un enunciado claro, tipo:
	“En esta sección presentamos hallazgos preliminares sobre la frecuencia y los contextos en los que aparece el término ciencia en las columnas de opinión analizadas (2018–2020).”
	
	2. Evidencia cuantitativa
	
	Gráficos o tablas simples:
	
	Tendencias temporales (ej. apariciones de ciencia por mes/año).
	
	Distribución por medio o autor.
	
	Top 20 términos que coocurren con ciencia.
	
	Ejemplo de redacción:
	“Entre 2018 y 2020, el término ciencia aparece en 2.340 columnas (29,2% del corpus). Su presencia es más recurrente en 2020, coincidiendo con la pandemia de COVID-19.”
	
	3. Evidencia cualitativa
	
	Fragmentos de columnas donde la palabra aparece con funciones diferentes (argumentativa, retórica, simbólica).
	
	Ejemplo:
	“En algunos casos, ciencia se emplea como fuente de legitimidad (‘la ciencia demuestra que…’), mientras que en otros se usa con valor metafórico (‘la ciencia política de la corrupción’).”
	
	4. Conexión con la pregunta de investigación
	
	Relaciona el hallazgo con las preguntas planteadas en el resumen:
	
	¿Qué usos del discurso científico identificamos hasta ahora?
	
	¿Cómo se valora la ciencia en la esfera pública?
	
	Ejemplo:
	“Estos hallazgos preliminares sugieren que la ciencia no solo se asocia a temas de salud y educación, sino que también funciona como metáfora en discusiones políticas, lo cual abre la pregunta sobre el grado en que la ciencia es apropiada simbólicamente en la opinión pública.”
	
	5. Visualización / síntesis
	
	Incluye una gráfica o tabla por hallazgo preliminar:
	
	Nube de palabras coocurrentes con ciencia.
	
	Gráfico de barras con los tópicos más frecuentes.
	
	Red semántica pequeña con actores y conceptos conectados.
	
	Ejemplo de redacción de un hallazgo preliminar:
	
	Hallazgo preliminar 1: Aparición y contextos del término ciencia
	
	Del total de 8.000 columnas analizadas, 2.340 (29,2%) incluyen al menos una mención explícita a ciencia. La frecuencia muestra un aumento pronunciado en 2020, en paralelo con la pandemia.
	
	En términos de coocurrencias, ciencia aparece asociada principalmente a salud, educación y tecnología. Sin embargo, también se emplea en contextos metafóricos vinculados a la política.
	
	Estos resultados preliminares sugieren que la ciencia cumple un doble papel en la esfera pública: como autoridad epistémica (ej. validación de argumentos) y como recurso retórico en debates no especializados.
	
	\section{Introducción}
	Motiva el problema, contribuciones y preguntas de investigación. Sitúa el trabajo en la literatura de comunicación pública de la ciencia, estudios CTS/STS y humanidades digitales (\citealp{BIBkey1,BIBkey2}).
	
	\section{Corpus y contexto}
	Describe fuentes (tres periódicos principales), rango temporal (2018--2020), criterios de inclusión/exclusión, metadatos (medio, autoría, fecha, URL). Discute representatividad y sesgos del corpus.
	
	\section{Metodología}
	\subsection{Preprocesamiento (PLN)}
	Limpieza, normalización, stopwords extendidas, lematización en español (spaCy/Stanza). 
	Define ventanas de contexto para analizar usos de “ciencia”.
	
	\subsection{Modelado y análisis}
	\begin{itemize}
		\item \textbf{Coocurrencias y redes semánticas:} PMI/PPMI; grafos conceptos–actores.
		\item \textbf{Tópicos:} BERTopic/NMF con validación por coherencia y revisión humana.
		\item \textbf{Valoración/stance:} esquema de anotación (acuerdo interanotador), clasificador ligero (LogReg/SVM) o modelo BETO/roberta-es; métricas (F1, AUC).
		\item \textbf{Dinámica temporal:} series mensuales y detección de puntos de cambio.
	\end{itemize}
	
	\section{Resultados preliminares}
	Resumen de hallazgos cuantitativos (frecuencias, tópicos, marcos) y cualitativos (funciones argumentativas, retóricas, simbólicas). 
	Refiere a figuras y tablas: ver \cref{fig:freq,tab:coocs}.
	
%	\begin{figure}[ht]
%		\centering
%		\includegraphics[width=0.8\linewidth]{figs/frecuencia_ciencia_por_mes.pdf}
%		\caption{Frecuencia mensual de menciones a \emph{ciencia} (2018--2020) por medio.}
%		\label{fig:freq}
%	\end{figure}
	
	\begin{table}[ht]
		\centering
		\caption{Top-15 coocurrencias de \emph{ciencia} (PPMI), normalizadas por 1{,}000 palabras.}
		\label{tab:coocs}
		\begin{tabular}{l S[table-format=1.3]}
			\toprule
			Término & {PPMI} \\
			\midrule
			salud & 2.341 \\
			educación & 2.115 \\
			tecnología & 1.982 \\
			% ... completa con tus resultados
			\bottomrule
		\end{tabular}
	\end{table}
	
	\section{Discusión}
	Interpreta los patrones: autoridad epistémica vs. recurso retórico; entrelazamientos con política/ética/educación. Conecta con literatura CTS/STS y comunicación pública de la ciencia.
	
	\section{Limitaciones y consideraciones éticas}
	Cobertura mediática, sesgos de autoría/medio, ambigüedad semántica, cumplimiento de TOS/licencias, privacidad y reproducibilidad.
	
	\section{Reproducibilidad y datos}
	Enlace a repositorio (GitHub/OSF/Zenodo), scripts, versión de paquetes y \emph{conteiner} (\texttt{Dockerfile}) si aplica. 
	Describe estructura del repo y cómo reproducir figuras/tablas.
	
	\section{Conclusiones y trabajo futuro}
	Síntesis de aportes y próximos pasos (ampliar periodos, añadir encuestas, enriquecer anotaciones, análisis multilingüe, etc.).
	
	\section*{Agradecimientos}
	Reconocimientos a colaboradores, financiación, y fuentes de datos.
	
	% --- Bibliografía ---
	\bibliography{references}  % Crea un archivo references.bib con tus entradas
	
	% --- Apéndices (opcional) ---
	\appendix
	\section{Guía mínima para citas y figuras}
	Ejemplo de cita: \citep{BIBkey1}; cita narrativa: \citet{BIBkey2}.
	Ecuación de ejemplo:
	\begin{equation}
	\text{PPMI}(w,c) = \max\left\{0, \log \frac{p(w,c)}{p(w)\,p(c)}\right\}.
	\end{equation}
	
\end{document}
